%% kapa: exam=1 round=19 year=2008 period=1 form=A subject=appraisal_law subject_ko=부동산관계법규 seq=7 number=47
%% source: Q-Net 감정평가사 기출문제 (cst00309) artlSeq=1001962 file=extracted/1008/1001962/2008년 제19회 감정평가사 1차시험 기출문제/1교시.hwp
%% conversion: kordoc 4.24.51 -> tools/gen1.py (auto); review=proofread-sonnet; flags=none
\begin{question}{47}
도시관리계획의 입안제안제도에 관한 설명으로 옳은 것은?
\begin{choices}
  \item 도시관리계획의 입안제안권자는 주민이며, 이해관계자는 여기에서의 주민에 포함되지 아니한다.
  \item 주민은 도시관리계획 입안권자에게 기반시설의 설치․정비․개량 및 용도지역의 변경에 관한 사항에 대하여 도시관리계획의 입안을 제안할 수 있다.
  \item 도시관리계획 입안제안서에는 도시관리계획도서, 계획설명서 및 문화재지표조사설명서를 첨부하여야 한다.
  \item 주민은 지구단위계획구역의 지정․변경사항에 대해서는 도시관리계획 입안권자에게 그 입안을 제안할 수 없다.
  \item 도시관리계획입안의 제안을 받은 도시관리계획 입안권자는 그 제안일부터 60일 이내에 도시관리계획입안에의 반영 여부를 제안자에게 통보하여야 하며, 부득이한 사정이 있는 경우에는 1회에 한하여 30일을 연장할 수 있다.
\end{choices}
\end{question}
\answer{5}
